% Options for packages loaded elsewhere
\PassOptionsToPackage{unicode}{hyperref}
\PassOptionsToPackage{hyphens}{url}
\documentclass[
]{article}
\usepackage{xcolor}
\usepackage{amsmath,amssymb}
\setcounter{secnumdepth}{-\maxdimen} % remove section numbering
\usepackage{iftex}
\ifPDFTeX
  \usepackage[T1]{fontenc}
  \usepackage[utf8]{inputenc}
  \usepackage{textcomp} % provide euro and other symbols
\else % if luatex or xetex
  \usepackage{unicode-math} % this also loads fontspec
  \defaultfontfeatures{Scale=MatchLowercase}
  \defaultfontfeatures[\rmfamily]{Ligatures=TeX,Scale=1}
\fi
\usepackage{lmodern}
\ifPDFTeX\else
  % xetex/luatex font selection
\fi
% Use upquote if available, for straight quotes in verbatim environments
\IfFileExists{upquote.sty}{\usepackage{upquote}}{}
\IfFileExists{microtype.sty}{% use microtype if available
  \usepackage[]{microtype}
  \UseMicrotypeSet[protrusion]{basicmath} % disable protrusion for tt fonts
}{}
\makeatletter
\@ifundefined{KOMAClassName}{% if non-KOMA class
  \IfFileExists{parskip.sty}{%
    \usepackage{parskip}
  }{% else
    \setlength{\parindent}{0pt}
    \setlength{\parskip}{6pt plus 2pt minus 1pt}}
}{% if KOMA class
  \KOMAoptions{parskip=half}}
\makeatother
\usepackage{longtable,booktabs,array}
\newcounter{none} % for unnumbered tables
\usepackage{calc} % for calculating minipage widths
% Correct order of tables after \paragraph or \subparagraph
\usepackage{etoolbox}
\makeatletter
\patchcmd\longtable{\par}{\if@noskipsec\mbox{}\fi\par}{}{}
\makeatother
% Allow footnotes in longtable head/foot
\IfFileExists{footnotehyper.sty}{\usepackage{footnotehyper}}{\usepackage{footnote}}
\makesavenoteenv{longtable}
\setlength{\emergencystretch}{3em} % prevent overfull lines
\providecommand{\tightlist}{%
  \setlength{\itemsep}{0pt}\setlength{\parskip}{0pt}}
\usepackage{bookmark}
\IfFileExists{xurl.sty}{\usepackage{xurl}}{} % add URL line breaks if available
\urlstyle{same}
\hypersetup{
  hidelinks,
  pdfcreator={LaTeX via pandoc}}

\author{}
\date{}

\begin{document}

\section{An Observation on the Structural Correspondences between
Signal/Control Theory and Quantum Mechanics, Quantum Optics, and Open
Quantum
Systems}\label{an-observation-on-the-structural-correspondences-between-signalcontrol-theory-and-quantum-mechanics-quantum-optics-and-open-quantum-systems}

\textbf{Author}: Noriaki Kihara \textbf{Affiliation}: WF System Co.,
Ltd.~/ Faculty of Engineering Science, Osaka University (graduate)
\textbf{ORCID}:
\href{https://orcid.org/0009-0004-6753-4020}{0009-0004-6753-4020}
\textbf{Version}: v6 \textbf{Date}: 3 June 2026 \textbf{License}: CC BY
4.0

\begin{center}\rule{0.5\linewidth}{0.5pt}\end{center}

\subsection{Character of This Note}\label{character-of-this-note}

\textbf{This is an observation paper. It is neither a proof paper nor an
assertion paper.}

This note makes no new physical prediction. It proves no new
mathematical theorem. It proposes no new interpretation.

What it does is limited to \textbf{juxtaposing and observing that ten
structural correspondences already exist in the literature} between the
mathematical structures already established and socially implemented in
signal/control theory and the mathematical structures established in
quantum mechanics, quantum optics, and open quantum systems.

These correspondences are facts already recorded in the standard
textbooks and classical original papers of both fields. This note merely
organizes them across fields.

This note does \textbf{not} claim:

\begin{itemize}
\tightlist
\item
  to propose a new interpretation of quantum mechanics
\item
  to claim a physical extension of signal/control theory
\item
  to claim the superiority of either field
\item
  to predict physical constants
\item
  to propose a modification of existing physical theory
\item
  to assert new mathematical theorems
\end{itemize}

Evaluation and interpretation are left to the reader.

\begin{center}\rule{0.5\linewidth}{0.5pt}\end{center}

\subsection{Abstract}\label{abstract}

Theoretical physics (quantum mechanics, quantum optics, open quantum
systems) and engineering (signal/control theory) developed independently
in the 20th century. The former aimed to describe microscopic physical
phenomena; the latter aimed at the practical realization of
communication, measurement, and control.

Yet the mathematical structures underlying the two fields have, at
several important points, an exact identity or a strong structural
correspondence. In particular, Heisenberg's quantum-mechanical
uncertainty \(\Delta x \cdot \Delta p \geq \hbar/2\) (Heisenberg 1927)
and Gabor's time--band uncertainty
\(\Delta t \cdot \Delta \omega \geq 1/2\) (Gabor 1946) are the same
mathematical theorem, and Gabor himself made the equivalence explicit.

Starting from this known correspondence, this note organizes the
structural correspondences of the two fields into ten items. For each
item, the standard literature of both fields is cited side by side, and
the range in which an exact identity holds and the range that remains a
structural parallel are made explicit.

As noteworthy facts, we observe: (i) the same mathematical object was
discovered independently in the two fields (Wigner 1932 in quantum
mechanics, Ville 1948 in signal theory, the equivalent phase-space
quasi-probability distribution); (ii) the Kennard--Robertson-type
Fourier uncertainty and Gabor's time--band uncertainty are based on the
same Fourier-analytic inequality, and Gabor himself made the
correspondence explicit; (iii) the mathematical structures on the
signal/control side are already socially implemented and operating daily
in engineering domains such as communications, GPS, autonomous driving,
radar, MRI, and optical communications.

\begin{center}\rule{0.5\linewidth}{0.5pt}\end{center}

\subsection{§1 Introduction}\label{introduction}

\subsubsection{1.1 Independent development of the two
fields}\label{independent-development-of-the-two-fields}

In the first half of the 20th century, quantum mechanics was
systematized by Heisenberg, Schrödinger, Dirac, von Neumann, and others
{[}1, 11, 16, 17{]}. In the same period, signal/control theory was
systematized by Nyquist, Shannon, Wiener, Gabor, Kalman, and others
{[}2, 5, 6, 7, 8{]}.

The two fields \textbf{developed independently with different aims,
motivations, and ranges of application}. The former aimed at the
prediction and understanding of microscopic physical phenomena; the
latter aimed at solving practical problems in communication,
measurement, and control.

\subsubsection{1.2 Correspondences of mathematical
structures}\label{correspondences-of-mathematical-structures}

Yet, juxtaposing the mathematical structures of the two fields, one can
confirm that \textbf{an exact identity or a strong structural
correspondence} holds for several important concepts. This note
organizes that observation.

The most explicit example is the identity between Heisenberg's
quantum-mechanical uncertainty principle {[}1{]} and Gabor's
signal-theoretic uncertainty principle {[}2{]}. Gabor (1946) {[}2{]}
explicitly pointed out in his own paper that the two have the same
mathematical structure.

\subsubsection{1.3 Method of this note}\label{method-of-this-note}

This note adopts the following method:

\begin{enumerate}
\def\labelenumi{\arabic{enumi}.}
\tightlist
\item
  Select pairs of concepts observed to have a correspondence from the
  two fields.
\item
  Cite, for each pair, the standard literature of both fields (classical
  original papers or canonical textbooks).
\item
  Juxtapose the correspondences of mathematical structure in table form.
\item
  \textbf{Make explicit the range in which an exact identity holds and
  the range that remains a structural parallel.}
\item
  Add no interpretation or evaluation.
\end{enumerate}

This note is limited to organizing facts obvious to experts in both
fields; its novelty lies only in the ``juxtaposition and
classification'' of structure. It makes no new mathematical or physical
claim.

\subsubsection{1.4 Structure of this note}\label{structure-of-this-note}

§2 lists the ten structural correspondences in turn. For each item it
distinguishes ``the exactly identical part'' from ``the part remaining a
structural parallel.'' §3 summarizes the observations. §4 makes explicit
what this note does not claim. §5 lists implications (these too are
observations, not claims). §6 gives the conclusion.

\begin{center}\rule{0.5\linewidth}{0.5pt}\end{center}

\subsection{§2 Ten Structural
Correspondences}\label{ten-structural-correspondences}

\subsubsection{2.0 Classification of the
correspondences}\label{classification-of-the-correspondences}

We classify the following ten items by the strength of correspondence.

{\def\LTcaptype{none} % do not increment counter
\begin{longtable}[]{@{}
  >{\raggedright\arraybackslash}p{(\linewidth - 4\tabcolsep) * \real{0.3333}}
  >{\raggedright\arraybackslash}p{(\linewidth - 4\tabcolsep) * \real{0.3333}}
  >{\raggedright\arraybackslash}p{(\linewidth - 4\tabcolsep) * \real{0.3333}}@{}}
\toprule\noalign{}
\begin{minipage}[b]{\linewidth}\raggedright
\#
\end{minipage} & \begin{minipage}[b]{\linewidth}\raggedright
Concept pair
\end{minipage} & \begin{minipage}[b]{\linewidth}\raggedright
Classification
\end{minipage} \\
\midrule\noalign{}
\endhead
\bottomrule\noalign{}
\endlastfoot
1 & Uncertainty principle & \textbf{Exact mathematical identity} (same
theorem up to unit scale) \\
2 & Time--frequency quasi-probability distribution & \textbf{Exact
mathematical identity} (identical defining formula) \\
3 & Paraxial equation / Schrödinger equation & \textbf{Conditional exact
isomorphism} (identical as a PDE form under the paraxial, monochromatic,
scalar approximation) \\
4 & Sampling / phase-space degrees of freedom & \textbf{Strong
structural correspondence} (formal correspondence of effective-DOF
counting in a finite domain; asymptotic identity of effective DOF) \\
5 & State-space representation / Hilbert-space picture &
\textbf{Structural parallel} (linear time evolution, but differences in
unitarity and input/output structure) \\
6 & Observability / CSCO & \textbf{Structural parallel} (parallel formal
problems, but mechanistic difference of dynamics/kinematics) \\
7 & Kalman filter / quantum filtering & \textbf{Structural parallel}
(corresponds via POVM, distinct from projective measurement) \\
8 & Jones vector / qubit & \textbf{Exact mathematical identity} (same
two-component complex vector, \(U(2)\) action) \\
9 & Dephasing / phase noise & \textbf{Structural parallel} (shares only
the mathematics of the \(T_2\) process, not decoherence as a whole) \\
10 & SVD / Schmidt decomposition & \textbf{Exact mathematical identity}
(same linear-algebra tool) \\
\end{longtable}
}

Exact mathematical identity holds for the 5 items \#1, \#2, \#3, \#8,
\#10; strong structural correspondence for the 1 item \#4; and
structural parallel for the 4 items \#5, \#6, \#7, \#9.

\subsubsection{2.1 Uncertainty principle (exact mathematical
identity)}\label{uncertainty-principle-exact-mathematical-identity}

{\def\LTcaptype{none} % do not increment counter
\begin{longtable}[]{@{}
  >{\raggedright\arraybackslash}p{(\linewidth - 2\tabcolsep) * \real{0.5000}}
  >{\raggedright\arraybackslash}p{(\linewidth - 2\tabcolsep) * \real{0.5000}}@{}}
\toprule\noalign{}
\begin{minipage}[b]{\linewidth}\raggedright
Signal theory
\end{minipage} & \begin{minipage}[b]{\linewidth}\raggedright
Quantum mechanics
\end{minipage} \\
\midrule\noalign{}
\endhead
\bottomrule\noalign{}
\endlastfoot
\textbf{Gabor time--band uncertainty} &
\textbf{Heisenberg--Kennard--Robertson uncertainty} \\
\(\Delta t \cdot \Delta \omega \geq \dfrac{1}{2}\) &
\(\Delta x \cdot \Delta p \geq \dfrac{\hbar}{2}\) \\
{[}Gabor 1946{]} {[}2{]} & {[}Heisenberg 1927{]} {[}1{]}, {[}Kennard
1927{]} {[}22{]}, {[}Robertson 1929{]} {[}23{]} \\
\end{longtable}
}

The two are the same inequality up to units (natural units \(\hbar=1\)).
The mathematical proofs of both rest on the Cauchy--Schwarz inequality
and the unitarity of the Fourier transform, and are structurally
identical. Since the Fourier conjugate of the position-\(x\)
representation is the momentum \(p/\hbar=k\),
\(\sigma_x\sigma_p\geq\hbar/2\) is the same Fourier-analytic proposition
as the Gabor inequality.

\textbf{Historical note}: Heisenberg (1927) {[}1{]} presented the
uncertainty principle as an intuitive, operational argument. It was
Kennard (1927) {[}22{]} who formulated it as the modern strict variance
inequality \(\sigma_x\sigma_p\geq\hbar/2\), and Robertson (1929)
{[}23{]} who generalized it to an arbitrary pair of self-adjoint
operators. The ``Heisenberg uncertainty'' referred to in this note means
the modern standard form based on these three papers.

Gabor (1946) {[}2{]} explicitly stated in his own work that his
inequality is ``an analog of Heisenberg's uncertainty principle in wave
mechanics.'' That is, \textbf{the equivalence was already stated, as of
1946, by the founder on the signal-theory side himself.}

\subsubsection{2.2 Time--frequency (phase-space) quasi-probability
distribution (exact mathematical
identity)}\label{timefrequency-phase-space-quasi-probability-distribution-exact-mathematical-identity}

{\def\LTcaptype{none} % do not increment counter
\begin{longtable}[]{@{}
  >{\raggedright\arraybackslash}p{(\linewidth - 2\tabcolsep) * \real{0.5000}}
  >{\raggedright\arraybackslash}p{(\linewidth - 2\tabcolsep) * \real{0.5000}}@{}}
\toprule\noalign{}
\begin{minipage}[b]{\linewidth}\raggedright
Signal theory
\end{minipage} & \begin{minipage}[b]{\linewidth}\raggedright
Quantum mechanics
\end{minipage} \\
\midrule\noalign{}
\endhead
\bottomrule\noalign{}
\endlastfoot
\textbf{Wigner--Ville distribution} & \textbf{Wigner function} \\
\(W(t,\omega) = \int s(t+\tau/2)\,\overline{s(t-\tau/2)}\, e^{-i\omega\tau}\, d\tau\)
&
\(W(x,p) = \dfrac{1}{2\pi\hbar}\int \psi^*(x+y/2)\,\psi(x-y/2)\, e^{ipy/\hbar}\, dy\) \\
{[}Ville 1948{]} {[}4{]} & {[}Wigner 1932{]} {[}3{]} \\
\end{longtable}
}

Both are quasi-probability distributions following the same mathematical
definition. \textbf{Wigner (quantum mechanics, 1932 {[}3{]}) and Ville
(signal theory, 1948 {[}4{]}) are different people}, yet the
mathematical objects they discovered independently have the same
structure. That the standard name in both fields, the
``\textbf{Wigner--Ville distribution},'' carries both names reflects the
history of independent discovery and mathematical identity.

Both distributions have the same properties:

\begin{itemize}
\tightlist
\item
  their marginals give the true probability distributions
  (position/momentum or time/frequency)
\item
  they can take negative values (quasi-probability, not satisfying
  positivity)
\item
  covariance under the linear Fourier transform
\item
  the same autocorrelation structure
\end{itemize}

Both are used daily, \textbf{as the same mathematical object}, in signal
processing (radar, sonar, time--frequency analysis) and in quantum
optics (quantum tomography).

\subsubsection{2.3 Schrödinger equation and paraxial wave equation
(conditional exact isomorphism: PDE form under the paraxial
approximation)}\label{schruxf6dinger-equation-and-paraxial-wave-equation-conditional-exact-isomorphism-pde-form-under-the-paraxial-approximation}

{\def\LTcaptype{none} % do not increment counter
\begin{longtable}[]{@{}
  >{\raggedright\arraybackslash}p{(\linewidth - 2\tabcolsep) * \real{0.5000}}
  >{\raggedright\arraybackslash}p{(\linewidth - 2\tabcolsep) * \real{0.5000}}@{}}
\toprule\noalign{}
\begin{minipage}[b]{\linewidth}\raggedright
Signal theory (optics)
\end{minipage} & \begin{minipage}[b]{\linewidth}\raggedright
Quantum mechanics
\end{minipage} \\
\midrule\noalign{}
\endhead
\bottomrule\noalign{}
\endlastfoot
\textbf{Paraxial wave equation} & \textbf{Schrödinger equation} \\
\(i\dfrac{\partial E}{\partial z} = -\dfrac{1}{2k}\nabla_\perp^2 E - \dfrac{k\,\delta n}{n_0}E\)
&
\(i\hbar \dfrac{\partial \psi}{\partial t} = -\dfrac{\hbar^2}{2m}\nabla^2 \psi + V\psi\) \\
Optical fiber, laser propagation, holography & Single-particle
non-relativistic quantum mechanics \\
{[}Saleh \& Teich 1991{]} {[}12{]} & {[}Schrödinger 1926{]} {[}16{]} \\
\end{longtable}
}

Under the \textbf{paraxial, monochromatic, scalar approximation}, the
two become \textbf{mathematically isomorphic partial differential
equations} under the correspondence of variables
\((z \leftrightarrow t,\ 1/k \leftrightarrow \hbar/m,\ V/\hbar \leftrightarrow -\,k\delta n/n_0)\).
Here, adopting the standard carrier convention \(E\propto e^{ikz}\), the
refractive-index term appears with a negative sign, and a high-index
region (\(\delta n>0\), waveguiding) corresponds to a \textbf{bound
potential well (\(V<0\))} (relative minus sign). Note that a repulsive
potential (\(V>0\)) is also a legitimate Schrödinger form; this sign
does not affect the isomorphism classification and is a note for the
consistency of the explicit dictionary.

This isomorphism is used daily in the design of optical-fiber
communications and integrated optics, and is made explicit in standard
optics textbooks including Saleh \& Teich (1991) {[}12{]}, Chapters
2--3. On the optics side it is sometimes taught as ``borrowing the
mathematics of quantum mechanics for optics,'' and sometimes the
reverse, ``borrowing the mathematics of optics for quantum mechanics'';
in either direction the mathematical content is the same.

\textbf{Caution}: the correspondence is, after all, an approximate,
mathematical isomorphism. On the optics side \(z\) is a propagation
distance; on the quantum side \(t\) is time. The physical meaning
(probability interpretation, Hilbert-space inner-product structure,
measurement theory, etc.) is not identical; the correspondence stays at
the mathematical form of a linear Schrödinger-type PDE.

\subsubsection{2.4 Sampling theorem and phase-space degrees of freedom
(strong structural
correspondence)}\label{sampling-theorem-and-phase-space-degrees-of-freedom-strong-structural-correspondence}

{\def\LTcaptype{none} % do not increment counter
\begin{longtable}[]{@{}
  >{\raggedright\arraybackslash}p{(\linewidth - 2\tabcolsep) * \real{0.5000}}
  >{\raggedright\arraybackslash}p{(\linewidth - 2\tabcolsep) * \real{0.5000}}@{}}
\toprule\noalign{}
\begin{minipage}[b]{\linewidth}\raggedright
Signal theory
\end{minipage} & \begin{minipage}[b]{\linewidth}\raggedright
Quantum mechanics
\end{minipage} \\
\midrule\noalign{}
\endhead
\bottomrule\noalign{}
\endlastfoot
\textbf{Shannon--Nyquist sampling theorem / time-bandwidth product} &
\textbf{Phase-space degree-of-freedom counting} \\
A real signal of band \(B\) and observation time \(T\) is approximately
described by an effective DOF \(N\sim 2BT\) (for complex baseband,
\(\sim BT\) independent \textbf{complex} samples) & The number of
quantum states for a phase-space volume \(V\cdot p_{\max}^d\) is
\(\sim V\cdot p_{\max}^d/(2\pi\hbar)^d\) \\
{[}Nyquist 1928{]} {[}5{]}, {[}Shannon 1949{]} {[}6{]},
{[}Slepian--Pollak 1961--1964{]} {[}13{]} & {[}Planck 1906{]} {[}24{]},
{[}Sackur 1911; Tetrode 1912{]} {[}25{]}, standard statistical-mechanics
textbooks \\
\end{longtable}
}

The two have a \textbf{strong structural correspondence} based on the
common framework of the concept of a ``\textbf{minimal phase-space
cell}.'' In signal theory it appears as the time-bandwidth product; in
quantum mechanics as the phase-space volume divided by the Planck unit
\(h\).

\textbf{The form of the DOF count corresponds strongly}, but the two are
not strictly the same theorem. The Shannon--Nyquist /
Slepian--Pollak-type time--band DOF and the semiclassical number of
states (the quantum number of states by the Weyl law) each have a
different measure, boundary condition, and physical interpretation.
Gabor (1946) {[}2{]} introduced this DOF count on the signal-theory side
under the name ``\textbf{logon},'' a counting essentially parallel to
the phase-space-cell concept on the quantum side.

\textbf{Refinement}: on the signal-theory side, the strict formulation
of the effective DOF in a finite-time, finite-band signal space is
established by the prolate spheroidal wave functions of Slepian, Pollak,
and Landau {[}13{]}. \textbf{Since a nonzero signal that is both
time-limited and band-limited does not strictly exist, the signal space
is formally infinite-dimensional, and \(N\sim 2BT\) is the asymptotic
expression of ``the effective dimension where the concentrated (prolate)
eigenvalues are close to 1.''} On the quantum side, the count dividing
the phase-space volume by \(h^d\) originates in Planck's quantization of
blackbody radiation (1906) {[}24{]} and the Sackur--Tetrode equation of
statistical mechanics (Sackur 1911; Tetrode 1912) {[}25{]}, and is
generalized as the Weyl law for the semiclassical number of states.

Stated information-theoretically, ``\textbf{the independent DOF of a
finite-band, finite-observation-time system is effectively finite},''
which is the same kind of framework as the finitization of the
Hilbert-space dimension in quantum mechanics (by phase-space volume).
The correspondence of the two is a strong structural correspondence of
\textbf{the same form of DOF counting, not the same theorem}.

\subsubsection{2.5 State-space representation and Hilbert-space picture
(structural parallel, with important
differences)}\label{state-space-representation-and-hilbert-space-picture-structural-parallel-with-important-differences}

{\def\LTcaptype{none} % do not increment counter
\begin{longtable}[]{@{}
  >{\raggedright\arraybackslash}p{(\linewidth - 2\tabcolsep) * \real{0.5000}}
  >{\raggedright\arraybackslash}p{(\linewidth - 2\tabcolsep) * \real{0.5000}}@{}}
\toprule\noalign{}
\begin{minipage}[b]{\linewidth}\raggedright
Control theory
\end{minipage} & \begin{minipage}[b]{\linewidth}\raggedright
Quantum mechanics
\end{minipage} \\
\midrule\noalign{}
\endhead
\bottomrule\noalign{}
\endlastfoot
\textbf{Linear state equation} & \textbf{Schrödinger equation (linear
version)} \\
State eq.: \(\dot{\mathbf{x}}(t) = A\,\mathbf{x}(t) + B\,\mathbf{u}(t)\)
& State evolution:
\(i\hbar\,\dfrac{d|\psi\rangle}{dt} = \hat{H}\,|\psi\rangle\) \\
Output eq.: \(\mathbf{y}(t) = C\,\mathbf{x}(t)\) & Observed value:
\(\langle A\rangle = \langle\psi|\hat{A}|\psi\rangle\) \\
{[}Kalman 1960{]} {[}8{]} & {[}von Neumann 1932{]} {[}11{]} \\
\end{longtable}
}

Both are parallel in that \textbf{a linear operator (\(A\) or
\(\hat{H}\)) describes a time evolution on a linear space}. However, the
following important differences exist:

\textbf{Difference (i): presence/absence of unitarity}. In quantum
mechanics \(\hat{H}\) is Hermitian (self-adjoint), so
\(e^{-i\hat{H}t/\hbar}\) is a unitary operator and the norm
(probability) of the state is conserved. By contrast, the \(A\) of
control theory need not be skew-Hermitian (purely imaginary spectrum) in
general, and \textbf{allows eigenvalues with a real part (unstable
poles, dissipative modes)}. That is, the conservative time evolution of
a closed quantum system and a control system handling open/dissipative
systems correspond only in the special case where \(A\) is
skew-Hermitian and \(B\mathbf{u}=0\).

\textbf{Difference (ii): bilinearity of the observed value}. The quantum
expectation \(\langle A\rangle=\langle\psi|\hat{A}|\psi\rangle\) is
\textbf{bilinear (a quadratic form)} in the state \(|\psi\rangle\),
whereas the control-theory output \(\mathbf{y}=C\mathbf{x}\) is
\textbf{linear} in the state \(\mathbf{x}\). The way observables are
extracted differs structurally.

\textbf{Difference (iii): presence/absence of a driving input}. The
driving input \(B\mathbf{u}(t)\) of control theory is the essence of a
controlled system, but the closed Schrödinger equation has no free term
corresponding to it. Extending to open quantum systems (Lindblad
equation {[}15{]}) adds a driving term, but this correspondence is
restricted to the standard closed Schrödinger equation.

Hence the correspondence of the two is a structural parallel at the
level of ``linear time evolution on a linear space,'' and even
restricted to finite-dimensional quantum systems (qubit, qudit),
\textbf{essential differences remain in unitarity, the way observables
are extracted, and the treatment of the driving input}. These
differences are made explicit in the literature of both fields and
indicate the limits of the structural correspondence.

\subsubsection{2.6 Observability and state distinguishability
(structural parallel, with mechanistic
difference)}\label{observability-and-state-distinguishability-structural-parallel-with-mechanistic-difference}

{\def\LTcaptype{none} % do not increment counter
\begin{longtable}[]{@{}
  >{\raggedright\arraybackslash}p{(\linewidth - 2\tabcolsep) * \real{0.5000}}
  >{\raggedright\arraybackslash}p{(\linewidth - 2\tabcolsep) * \real{0.5000}}@{}}
\toprule\noalign{}
\begin{minipage}[b]{\linewidth}\raggedright
Control theory
\end{minipage} & \begin{minipage}[b]{\linewidth}\raggedright
Quantum mechanics
\end{minipage} \\
\midrule\noalign{}
\endhead
\bottomrule\noalign{}
\endlastfoot
\textbf{Observability} & \textbf{Complete Set of Commuting Observables
(CSCO)} \\
Rank condition of the observability matrix
\(\mathcal{O} = \begin{pmatrix}C\\CA\\CA^2\\\vdots\\CA^{n-1}\end{pmatrix}\)
& The maximal set of mutually commuting self-adjoint operators \\
Whether the state is recoverable from the output time series & Whether
the state is specifiable from the simultaneous eigenvalues of
observables \\
{[}Kalman 1960{]} {[}8{]} & {[}Dirac 1930{]} {[}17{]}, {[}von Neumann
1932{]} {[}11{]} \\
\end{longtable}
}

The two are \textbf{parallel formalizations} of ``the condition for
specifying the state from observation.'' But there is an important
mechanistic difference:

\textbf{Difference (i): dynamics vs kinematics}. Kalman's observability
condition is the condition for uniquely determining the initial state
\(\mathbf{x}(0)\) from the output time series \(\{y(t)\}_{t\geq0}\),
using the \textbf{dynamics} of stacking \(CA^k\). By contrast, the
quantum CSCO is the \textbf{kinematic} condition that a commuting
observable algebra at a single time uniquely specifies a basis of the
state space, requiring no time evolution.

\textbf{Difference (ii): single trajectory vs many copies}. Kalman
observability reconstructs the state from a deterministic single
trajectory. In quantum mechanics, by contrast, a single measurement
generally cannot determine the state, and state determination requires
\textbf{quantum state tomography} (repeated measurement on many copies
of the same state), because quantum measurement is inherently
probabilistic and measurement disturbs the state.

The correspondence is a meta-level parallel of ``the criterion for
whether the observable algebra separates the state space,'' and the
concrete content of the mechanism differs in the two fields.

\textbf{Note on the correspondence}: From the viewpoint of state
recoverability, the quantum counterpart closest to control-theory
observability is not the CSCO but quantum state tomography or an
informationally complete POVM. The correspondence with the CSCO made in
this note is restricted to the kinematic aspect that observables
separate the state space.

\subsubsection{2.7 Optimal estimation and filtering (structural
parallel, correspondence via
POVM)}\label{optimal-estimation-and-filtering-structural-parallel-correspondence-via-povm}

{\def\LTcaptype{none} % do not increment counter
\begin{longtable}[]{@{}
  >{\raggedright\arraybackslash}p{(\linewidth - 2\tabcolsep) * \real{0.5000}}
  >{\raggedright\arraybackslash}p{(\linewidth - 2\tabcolsep) * \real{0.5000}}@{}}
\toprule\noalign{}
\begin{minipage}[b]{\linewidth}\raggedright
Control theory
\end{minipage} & \begin{minipage}[b]{\linewidth}\raggedright
Quantum mechanics
\end{minipage} \\
\midrule\noalign{}
\endhead
\bottomrule\noalign{}
\endlastfoot
\textbf{Kalman filter} & \textbf{Quantum filtering} \\
Predict → observe → update cycle & State update by generalized
measurement (POVM) and quantum conditional expectation \\
Minimum-mean-square-error estimation for a linear Gaussian state-space
model & Bayes-optimal estimation of the quantum state under classical
output measurement \\
{[}Kalman 1960{]} {[}8{]} & {[}Davies \& Lewis 1970{]} {[}18{]},
{[}Belavkin 1992{]} {[}19{]} \\
\end{longtable}
}

Both are parallel in being operations that ``\textbf{optimally update
the state using observed information}.'' But the center of gravity of
the correspondence lies not in \textbf{projective measurement (von
Neumann--Lüders) but in POVM and quantum filtering}:

\textbf{Important distinction}: projective measurement (von
Neumann--Lüders) involves an irreversible disturbing back-action on the
state, which has no counterpart in the Kalman filter. The Kalman filter
is a Bayes update on a classical probability distribution and shares no
mathematical structure with the state collapse of projective
measurement.

The correct bridge between the two passes through POVM (Positive
Operator-Valued Measure, Davies--Lewis 1970 {[}18{]}) and quantum
filtering theory (Belavkin 1992 {[}19{]}). Quantum filtering formulates,
in the framework of Bayes-optimal estimation, how the quantum state is
updated under continuous classical output measurement
(homodyne/heterodyne detection, etc.) (the Belavkin equation, the
stochastic master equation). \textbf{For linear Gaussian systems the
correspondence with a Kalman-type filter is strong, while general
quantum filtering remains a structural parallel of Bayesian state
update} (generalizing as ``the quantum version of the Kalman filter'' is
too strong).

That is, the correspondence of the Kalman filter with ``simple
projective measurement'' is inaccurate, and the correspondence of the
Kalman filter with ``POVM + quantum filtering'' is the mathematically
correct parallel.

\subsubsection{2.8 Jones vector and qubit (exact mathematical
identity)}\label{jones-vector-and-qubit-exact-mathematical-identity}

{\def\LTcaptype{none} % do not increment counter
\begin{longtable}[]{@{}
  >{\raggedright\arraybackslash}p{(\linewidth - 2\tabcolsep) * \real{0.5000}}
  >{\raggedright\arraybackslash}p{(\linewidth - 2\tabcolsep) * \real{0.5000}}@{}}
\toprule\noalign{}
\begin{minipage}[b]{\linewidth}\raggedright
Signal theory (optics/communications)
\end{minipage} & \begin{minipage}[b]{\linewidth}\raggedright
Quantum mechanics (optics)
\end{minipage} \\
\midrule\noalign{}
\endhead
\bottomrule\noalign{}
\endlastfoot
\textbf{Jones vector (polarization state)} & \textbf{qubit state / Bloch
sphere} \\
\(|\psi\rangle = \begin{pmatrix} E_x \\ E_y \end{pmatrix}\),
\(E_x, E_y \in \mathbb{C}\) &
\(|\psi\rangle = \alpha|0\rangle + \beta|1\rangle\),
\(\alpha, \beta \in \mathbb{C}\), \(|\alpha|^2 + |\beta|^2 = 1\) \\
Two-component complex vector (4 real DOF) & Two-component complex vector
(4 real DOF; 2 real DOF after normalization + global phase) \\
{[}Jones 1941{]} {[}9{]} & {[}Bloch 1946{]} {[}14{]}, standard
quantum-mechanics textbooks \\
\end{longtable}
}

Both represent the state by \textbf{the exact same two-component complex
vector}. Up to the normalization condition and the global-phase freedom,
both are described as points on the \textbf{complex projective line
\(\mathbb{CP}^1\cong S^2\)}, which has the same geometry as the Bloch
sphere (quantum mechanics) / Poincaré sphere (polarization).

Concrete correspondences:

\begin{itemize}
\tightlist
\item
  Orthogonal polarization-basis transformations (H/V ↔ ±45° ↔ R/L
  circular) ↔ qubit basis transformations (\(|0\rangle,|1\rangle\) vs
  \(|\pm\rangle\) vs \(|\pm i\rangle\))
\item
  \textbf{Lossless Jones matrices} (phase plates, rotators, birefringent
  elements, and other reversible polarization elements) ↔ \(U(2)\)
  \textbf{unitary transformations}
\item
  \textbf{General Jones matrices including absorption/polarizers} ↔
  \textbf{non-unitary linear transformations} (a formal correspondence
  with quantum operations including measurement/loss on the quantum
  side)
\item
  Stokes parameters ↔ Bloch-vector components
\item
  The \textbf{orthogonal polarization basis} used in
  polarization-division multiplexing (PDM) ↔ the \textbf{same
  polarization basis} used in quantum optics to describe single-photon
  polarization qubits and two-photon polarization states (caution: PDM
  is the dual-polarization multiplexing of classical optical
  communication, not quantum entanglement itself; what is shared is only
  the choice of polarization basis)
\end{itemize}

This is an exact mathematical identity restricted to pure polarization
states, and is used daily in the implementation of quantum optics and
quantum information. However, as above, the physical interpretation of
the general lossy Jones matrix and of PDM has no direct identity with
quantum unitarity or quantum entanglement.

\textbf{Caution}: this item treats the correspondence between the
\textbf{classical-optics Jones vector} and the \textbf{quantum qubit}.
The \textbf{I/Q complex baseband signal} \(\tilde{s}(t)=I(t)+iQ(t)\) in
communications engineering is a single complex scalar (2 real DOF), of
different dimension from the Jones vector (two-component complex, 4 real
DOF). What corresponds directly to the I/Q baseband signal is the
\textbf{single-mode complex amplitude} (the single-mode amplitude of a
classical electromagnetic field), not a qubit. Care must be taken not to
conflate the two.

\subsubsection{2.9 Dephasing and phase noise (structural parallel, only
part of
decoherence)}\label{dephasing-and-phase-noise-structural-parallel-only-part-of-decoherence}

{\def\LTcaptype{none} % do not increment counter
\begin{longtable}[]{@{}
  >{\raggedright\arraybackslash}p{(\linewidth - 2\tabcolsep) * \real{0.5000}}
  >{\raggedright\arraybackslash}p{(\linewidth - 2\tabcolsep) * \real{0.5000}}@{}}
\toprule\noalign{}
\begin{minipage}[b]{\linewidth}\raggedright
Signal theory
\end{minipage} & \begin{minipage}[b]{\linewidth}\raggedright
Quantum mechanics
\end{minipage} \\
\midrule\noalign{}
\endhead
\bottomrule\noalign{}
\endlastfoot
\textbf{Phase noise / jitter} & \textbf{Dephasing (pure phase
relaxation)} \\
\(\phi(t) = \omega_0 t + \delta\phi(t)\) & Dissipation of the
off-diagonal elements of the density matrix \\
Lorentzian linewidth (characteristic width \(\sim 1/T\),
convention-dependent) & Transverse relaxation time \(T_2\) \\
Oscillator theory, laser-linewidth theory ({[}Schawlow--Townes 1958{]}
{[}26{]}) & {[}Bloch 1946{]} {[}14{]}, {[}Lindblad 1976{]} {[}15{]} \\
\end{longtable}
}

Both share the mathematical structure of \textbf{dissipation of phase
information}. Concretely, it is described as a process in which the
phase distribution of the complex amplitude diffuses over time, and the
Lorentzian (the Schawlow--Townes linewidth in laser-linewidth theory
{[}26{]}) and the \(T_2\) relaxation of the Bloch equation have the same
mathematical form.

\textbf{On the linewidth convention}: the coefficient changes by
convention --- full width at half maximum (FWHM) or half width at half
maximum (HWHM), angular frequency \(\Delta\omega\) or ordinary frequency
\(\Delta f\) (e.g., for the Lorentzian spectrum of the exponential
correlation \(g^{(1)}(\tau)\propto e^{-|\tau|/T}\), FWHM \(=1/(\pi T)\)
Hz, HWHM \(=1/(2\pi T)\) Hz). This note states only the qualitative
isomorphism that the characteristic width is of order \(\sim 1/T\); the
exact coefficient correspondence depends on convention.

\textbf{Important restriction}: ``decoherence'' is a broad concept in
quantum information/optics and includes, in addition to (i) pure phase
relaxation (\(T_2\), dephasing), (ii) energy relaxation (\(T_1\),
transitions from excited to ground state) and (iii) loss of purity of
the reduced density matrix due to entanglement with the environment.
What strictly shares a mathematical structure with classical phase noise
is \textbf{only (i), pure phase relaxation (dephasing)}; (ii) and (iii)
have no classical counterpart (since classical systems have no discrete
energy levels or quantum entanglement).

Hence the correspondence of this item is restricted not to ``phase noise
= decoherence as a whole'' but to ``\textbf{phase noise = the dephasing
part of decoherence}.'' Within this restriction, \textbf{the
corresponding pure-dephasing process on the quantum side is described by
a Lindblad-form master equation {[}15{]}} (this means the quantum
counterpart is expressed in Lindblad form, not that classical phase
noise itself obeys a Lindblad equation).

\subsubsection{2.10 SVD and Schmidt decomposition (exact mathematical
identity)}\label{svd-and-schmidt-decomposition-exact-mathematical-identity}

{\def\LTcaptype{none} % do not increment counter
\begin{longtable}[]{@{}
  >{\raggedright\arraybackslash}p{(\linewidth - 2\tabcolsep) * \real{0.5000}}
  >{\raggedright\arraybackslash}p{(\linewidth - 2\tabcolsep) * \real{0.5000}}@{}}
\toprule\noalign{}
\begin{minipage}[b]{\linewidth}\raggedright
Signal theory (communications)
\end{minipage} & \begin{minipage}[b]{\linewidth}\raggedright
Quantum mechanics
\end{minipage} \\
\midrule\noalign{}
\endhead
\bottomrule\noalign{}
\endlastfoot
\textbf{Singular value decomposition (SVD)} & \textbf{Schmidt
decomposition} \\
\(H = U \Sigma V^\dagger\), \(H \in \mathbb{C}^{N_r \times N_t}\) &
\(|\Psi\rangle_{AB} = \sum_i \sqrt{\lambda_i} \, |a_i\rangle_A \otimes |b_i\rangle_B\) \\
Decompose an arbitrary matrix by orthogonal basis changes and singular
values & Decompose a bipartite pure state by Schmidt coefficients and
orthogonal bases \\
{[}Eckart--Young 1936{]} {[}20{]} & {[}Schmidt 1907{]} {[}21{]} \\
\end{longtable}
}

The two are \textbf{the same tool of linear algebra}. Representing a
bipartite pure state on the tensor-product space
\(\mathcal{H}_A\otimes\mathcal{H}_B\) as a coefficient matrix, its
singular value decomposition is the Schmidt decomposition itself. The
Schmidt coefficients \(\sqrt{\lambda_i}\) are the singular values of the
coefficient matrix, and the Schmidt bases
\(\{|a_i\rangle\},\{|b_i\rangle\}\) correspond to the left/right
orthogonal bases of the SVD. The historical origin of the concept of
bipartite entanglement in the quantum case goes back to
Einstein--Podolsky--Rosen (1935) {[}10{]}.

\textbf{Limits of the correspondence (important)}: what corresponds in
this item is only \textbf{the SVD = Schmidt decomposition as a tool of
linear algebra}. No further physical/engineering identity holds:

\begin{itemize}
\tightlist
\item
  \textbf{The MIMO channel matrix \(H\) and the coefficient matrix of a
  bipartite pure state are different objects}: the MIMO \(H\) is the
  matrix representing the propagation channel, and the signal vectors
  \(\mathbf{x},\mathbf{y}\) are the input/output vectors of a single
  linear channel. These are not ``vectors on a tensor-product space
  \(\mathcal{H}_A\otimes\mathcal{H}_B\) of two subsystems.'' By
  contrast, the quantum entangled state \(|\Psi\rangle_{AB}\) belongs to
  the tensor-product space of subsystems.
\item
  \textbf{The functional forms of the capacity formulas differ}: the
  MIMO capacity is \(C=\sum_i\log_2(1+\rho\sigma_i^2)\) (a logarithmic
  sum over singular values \(\sigma_i\) with SNR \(\rho\)), whereas the
  entanglement entropy is \(S=-\sum_i\lambda_i\log\lambda_i\) (the
  Shannon entropy of the squared, normalized Schmidt coefficients
  \(\lambda_i\)). The two have different functional forms, and the
  isomorphism of capacity does not hold.
\end{itemize}

Hence this item is limited to the observation that ``\textbf{the SVD, a
tool of linear algebra, plays the same role (giving the canonical form
of a bipartite system) in the two fields}.'' Since the structural
isomorphism between MIMO channel capacity and quantum entanglement
capacity does not hold, it is not claimed here.

\begin{center}\rule{0.5\linewidth}{0.5pt}\end{center}

\subsection{§3 Summary of the
Observations}\label{summary-of-the-observations}

All ten items above are facts made explicit in the standard textbooks
and classical original papers of both fields. This note merely
juxtaposes them across fields.

The following noteworthy facts can be observed:

\textbf{(A) Independent discovery by independent people}

A salient example of the same mathematical object being discovered
independently in the two fields is the phase-space quasi-probability
distribution. Wigner (1932 {[}3{]}) in quantum mechanics and Ville (1948
{[}4{]}) in signal theory introduced equivalent distributions
independently of each other. They are different people, and the standard
name in both fields, the ``\textbf{Wigner--Ville distribution},''
reflects their independent discovery. That the same mathematical object
was reached independently from different physical/engineering
motivations suggests the mathematical depth of the correspondence.

\textbf{(B) The same Fourier-analytic inequality, and Gabor's explicit
statement}

The Kennard--Robertson-type Fourier uncertainty and Gabor's time--band
uncertainty are based on the same Fourier-analytic inequality.
Importantly, Gabor himself made the mathematical identity of the two
explicit in his 1946 paper {[}2{]}. That is, the correspondence between
the two fields \textbf{was already recognized by the founders of the two
fields themselves.}

\textbf{(C) Social implementation in engineering}

The mathematical structures on the signal/control side are already
socially implemented and operating daily in the following engineering
domains:

\begin{itemize}
\tightlist
\item
  OFDM, MIMO, and I/Q demodulation in mobile communications (4G LTE, 5G
  NR, 6G research)
\item
  the Kalman filter in GPS and satellite positioning
\item
  state estimation in autonomous driving and drone control
\item
  the Wigner--Ville distribution in radar and sonar
\item
  numerical solution of the Schrödinger-type paraxial wave equation in
  optical-fiber communications
\item
  the Bloch equation in MRI (magnetic resonance imaging)
\item
  phase-noise analysis in oscillator and laser design
\end{itemize}

These implementations are a \textbf{confirmation of the operation of a
shared mathematical structure in engineering}. The empirical validity on
the quantum side is independently verified separately by quantum
experiments (quantum optics, atom interferometry, superconducting
qubits, etc.). The two stay at the observation that the same
mathematical structure is confirmed in different physical domains; the
social implementation of one does not directly guarantee the empirical
validity of the other (this distinction is important to avoid a category
confusion).

\begin{center}\rule{0.5\linewidth}{0.5pt}\end{center}

\subsection{§4 What This Note Does Not Claim (explicit
scope)}\label{what-this-note-does-not-claim-explicit-scope}

This note does not claim:

\begin{enumerate}
\def\labelenumi{\arabic{enumi}.}
\tightlist
\item
  \textbf{It does not propose a new interpretation of quantum
  mechanics}: the observed correspondences are, after all, a
  juxtaposition of mathematical structures, and express no view on the
  physical content or interpretation of quantum mechanics (Copenhagen,
  many-worlds, hidden variables, etc.).
\item
  \textbf{It does not claim a physical extension of signal/control
  theory}: signal/control theory is positioned purely as practical
  engineering, and no claim is made to re-position it as a new theory of
  physics.
\item
  \textbf{It does not claim the superiority of either field}: both
  fields developed independently and are both mature, so we are not in a
  position to judge superiority.
\item
  \textbf{It does not predict physical constants}: this note only
  juxtaposes structure; predicting constant values is out of scope.
\item
  \textbf{It does not propose a modification of existing physical
  theory}: it proposes no change to the mathematical structures of
  standard quantum mechanics, quantum optics, or open quantum systems.
\item
  \textbf{It does not assert new mathematical theorems}: all the
  structural correspondences are recorded in the existing literature of
  both fields. The contribution of this note lies only in juxtaposition
  and classification.
\item
  \textbf{It does not extend a strong identity to a structural
  parallel}: following the classification of §2.0, it distinguishes the
  range in which an exact mathematical identity holds from the range
  that remains a structural parallel, and makes no claim to extend a
  parallel to an identity.
\end{enumerate}

The evaluation of this note is left to the reader.

\begin{center}\rule{0.5\linewidth}{0.5pt}\end{center}

\subsection{§5 Implications (as
observations)}\label{implications-as-observations}

We record the following as implications that may be naturally drawn from
these observations. These are not claims but direct consequences of the
observations.

\textbf{(I) Consistency of the uncertainty principle}

The Gabor inequality established in signal/control theory and the
Heisenberg uncertainty principle of quantum mechanics are the same
mathematical structure (§2.1). \textbf{As long as one accepts the
content as a Fourier-analytic inequality, one cannot treat the
mathematical cores of the Gabor-type and Kennard-type uncertainties as
different things} (this does not claim an identity of physical meaning).

\textbf{(II) Operation of the shared mathematical structure in
engineering}

The mathematical structures on the signal/control side are already
socially implemented and operating daily in engineering domains such as
communications, GPS, autonomous driving, radar, and MRI (§3 (C)). This
means at least that part of the mathematical structure shared by the two
fields is empirically confirmed in engineering. The empirical validity
on the quantum side is confirmed separately by quantum experiments; the
two stay at the observation that the same mathematical structure is
independently confirmed in different physical domains.

\textbf{(III) Possibility of inter-language translation}

By cross-referencing the literature of the two fields, it is possible to
understand the concepts of each field in the language of the other
(especially for the exact-identity items of §2.0). For example,
Heisenberg uncertainty can be translated as the Gabor inequality, a
qubit as a Jones vector, and the Schrödinger equation as the paraxial
optics equation.

\textbf{(IV) Educational and communicative implications}

By making explicit the mathematical structures shared by theoretical
physics and engineering, it is observed that communication between
experts in the two fields, and bridging in interdisciplinary education,
become easier. However, for items classified as structural parallels in
§2.0, one must also understand the differences in mechanism and limits.

These implications too are not new claims but direct consequences of the
observations of §2.

\begin{center}\rule{0.5\linewidth}{0.5pt}\end{center}

\subsection{§6 Conclusion}\label{conclusion}

This note juxtaposed and observed, from the classical literature of both
fields, the \textbf{ten structural correspondences} already established
between signal/control theory and quantum mechanics, quantum optics, and
open quantum systems. Of the ten, 5 items (§2.1, §2.2, §2.3, §2.8,
§2.10) have an exact mathematical identity within the explicitly stated
conditions/conventions. §2.4 is not the same theorem but a
\textbf{strong structural correspondence} of effective-DOF counting in a
finite domain. The remaining 4 items (§2.5, §2.6, §2.7, §2.9) remain
structural parallels (see the classification table of §2.0).

These are not new claims but facts already made explicit in the standard
textbooks and original papers of both fields. The contribution of this
note lies in organizing, as a single list, the correspondences scattered
across the literature of both fields, and in distinguishing exact
identity from structural parallel.

The structural correspondences of the two fields have been recognized
gradually throughout the 20th century (Gabor 1946 {[}2{]} making the
Heisenberg equivalence explicit, the independent discovery by Wigner
1932 {[}3{]} and Ville 1948 {[}4{]}, the quantum-filtering theory since
Belavkin 1992 {[}19{]}, etc.). This note merely binds them into a single
observation.

We hope this note is useful as a bridging resource for researchers and
practitioners standing at the boundary of the two fields.

When readers a century hence reach a more comprehensive understanding of
the correspondences of the two fields, this note may be referred to as a
transitional organization. That is enough.

\begin{center}\rule{0.5\linewidth}{0.5pt}\end{center}

\subsection{References}\label{references}

{[}1{]} Heisenberg, W. (1927). ``Über den anschaulichen Inhalt der
quantentheoretischen Kinematik und Mechanik.'' \emph{Zeitschrift für
Physik}, \textbf{43}, 172--198. {[}2{]} Gabor, D. (1946). ``Theory of
Communication.'' \emph{Journal of the Institution of Electrical
Engineers -- Part III}, \textbf{93} (26), 429--457. {[}3{]} Wigner, E.
P. (1932). ``On the Quantum Correction for Thermodynamic Equilibrium.''
\emph{Physical Review}, \textbf{40}, 749--759. {[}4{]} Ville, J. (1948).
``Théorie et applications de la notion de signal analytique.''
\emph{Câbles et Transmission}, \textbf{2} (1), 61--74. {[}5{]} Nyquist,
H. (1928). ``Certain topics in telegraph transmission theory.''
\emph{Transactions of the AIEE}, \textbf{47} (2), 617--644. {[}6{]}
Shannon, C. E. (1949). ``Communication in the presence of noise.''
\emph{Proceedings of the IRE}, \textbf{37} (1), 10--21. {[}7{]} Wiener,
N. (1949). \emph{Extrapolation, Interpolation, and Smoothing of
Stationary Time Series}. MIT Press, Cambridge, MA. {[}8{]} Kalman, R. E.
(1960). ``A New Approach to Linear Filtering and Prediction Problems.''
\emph{Journal of Basic Engineering}, \textbf{82} (1), 35--45. {[}9{]}
Jones, R. C. (1941). ``A new calculus for the treatment of optical
systems.'' \emph{Journal of the Optical Society of America}, \textbf{31}
(7), 488--493. {[}10{]} Einstein, A., Podolsky, B., Rosen, N. (1935).
``Can Quantum-Mechanical Description of Physical Reality Be Considered
Complete?'' \emph{Physical Review}, \textbf{47}, 777--780. {[}11{]} von
Neumann, J. (1932). \emph{Mathematische Grundlagen der Quantenmechanik}.
Springer-Verlag, Berlin. {[}12{]} Saleh, B. E. A., Teich, M. C. (1991).
\emph{Fundamentals of Photonics}. Wiley-Interscience, New York. {[}13{]}
Slepian, D., Pollak, H. O. (1961). ``Prolate spheroidal wave functions,
Fourier analysis and uncertainty --- I.'' \emph{Bell System Technical
Journal}, \textbf{40} (1), 43--63; Landau, H. J., Pollak, H. O. (1961).
``\ldots{} --- II.'' \emph{Bell System Technical Journal}, \textbf{40}
(1), 65--84 (Parts III--V appeared in the same series, 1961--1978).
{[}14{]} Bloch, F. (1946). ``Nuclear Induction.'' \emph{Physical
Review}, \textbf{70} (7--8), 460--474. {[}15{]} Lindblad, G. (1976).
``On the generators of quantum dynamical semigroups.''
\emph{Communications in Mathematical Physics}, \textbf{48} (2),
119--130. {[}16{]} Schrödinger, E. (1926). ``Quantisierung als
Eigenwertproblem (Erste Mitteilung).'' \emph{Annalen der Physik},
\textbf{384} (4), 361--376. {[}17{]} Dirac, P. A. M. (1930). \emph{The
Principles of Quantum Mechanics}. Oxford University Press, Oxford.
{[}18{]} Davies, E. B., Lewis, J. T. (1970). ``An operational approach
to quantum probability.'' \emph{Communications in Mathematical Physics},
\textbf{17} (3), 239--260. {[}19{]} Belavkin, V. P. (1992). ``Quantum
stochastic calculus and quantum nonlinear filtering.'' \emph{Journal of
Multivariate Analysis}, \textbf{42} (2), 171--201. {[}20{]} Eckart, C.,
Young, G. (1936). ``The approximation of one matrix by another of lower
rank.'' \emph{Psychometrika}, \textbf{1} (3), 211--218. {[}21{]}
Schmidt, E. (1907). ``Zur Theorie der linearen und nichtlinearen
Integralgleichungen.'' \emph{Mathematische Annalen}, \textbf{63} (4),
433--476. {[}22{]} Kennard, E. H. (1927). ``Zur Quantenmechanik
einfacher Bewegungstypen.'' \emph{Zeitschrift für Physik}, \textbf{44},
326--352. {[}23{]} Robertson, H. P. (1929). ``The Uncertainty
Principle.'' \emph{Physical Review}, \textbf{34}, 163--164. {[}24{]}
Planck, M. (1906). \emph{Vorlesungen über die Theorie der
Wärmestrahlung}. Barth, Leipzig. {[}25{]} Sackur, O. (1911).
\emph{Annalen der Physik}, \textbf{341} (15), 958--980; Tetrode, H.
(1912). \emph{Annalen der Physik}, \textbf{343} (7), 434--442. {[}26{]}
Schawlow, A. L., Townes, C. H. (1958). ``Infrared and Optical Masers.''
\emph{Physical Review}, \textbf{112} (6), 1940--1949.

\begin{center}\rule{0.5\linewidth}{0.5pt}\end{center}

\subsection{Appendix A: Summary Table of
Correspondences}\label{appendix-a-summary-table-of-correspondences}

For the reader's convenience, the ten items of §2 are re-listed. The
\textbf{Classification} column distinguishes exact identity
(\textbf{{[}★ exact{]}}), strong structural correspondence (\textbf{{[}◎
strong{]}}), and structural parallel (\textbf{{[}△ parallel{]}}).

{\def\LTcaptype{none} % do not increment counter
\begin{longtable}[]{@{}
  >{\raggedright\arraybackslash}p{(\linewidth - 8\tabcolsep) * \real{0.2000}}
  >{\raggedright\arraybackslash}p{(\linewidth - 8\tabcolsep) * \real{0.2000}}
  >{\raggedright\arraybackslash}p{(\linewidth - 8\tabcolsep) * \real{0.2000}}
  >{\raggedright\arraybackslash}p{(\linewidth - 8\tabcolsep) * \real{0.2000}}
  >{\raggedright\arraybackslash}p{(\linewidth - 8\tabcolsep) * \real{0.2000}}@{}}
\toprule\noalign{}
\begin{minipage}[b]{\linewidth}\raggedright
\#
\end{minipage} & \begin{minipage}[b]{\linewidth}\raggedright
Signal/control theory
\end{minipage} & \begin{minipage}[b]{\linewidth}\raggedright
Quantum mechanics/optics/open systems
\end{minipage} & \begin{minipage}[b]{\linewidth}\raggedright
Common mathematical structure
\end{minipage} & \begin{minipage}[b]{\linewidth}\raggedright
Classification
\end{minipage} \\
\midrule\noalign{}
\endhead
\bottomrule\noalign{}
\endlastfoot
1 & Gabor time--band uncertainty & Heisenberg--Kennard--Robertson
uncertainty & Lower bound of the variance product of a Fourier conjugate
pair & \textbf{{[}★ exact{]}} \\
2 & Wigner--Ville time--frequency distribution & Wigner phase-space
quasi-probability & Fourier autocorrelation of a quadratic form (with
conventions/sign conventions) & \textbf{{[}★ exact{]}} \\
3 & Paraxial wave equation & Schrödinger equation & First-order-in-time,
second-order-in-space linear PDE (identical as a PDE form under the
paraxial/monochromatic/scalar approximation) & \textbf{{[}★ conditional
exact{]}} \\
4 & Shannon--Nyquist sampling & Phase-space DOF counting & Formal
correspondence of effective-DOF counting in a finite domain (asymptotic
identity of effective DOF) & \textbf{{[}◎ strong{]}} \\
5 & State/output equations & Schrödinger equation/expectation & Linear
time evolution on a linear space (differences in unitarity, etc.) &
\textbf{{[}△ parallel{]}} \\
6 & Kalman observability & CSCO & State separation by the observable
algebra (dynamics/kinematics difference) & \textbf{{[}△ parallel{]}} \\
7 & Kalman filter update & Quantum filtering (via POVM) & Bayes-optimal
state estimation under observation & \textbf{{[}△ parallel{]}} \\
8 & Jones vector (pure polarization) & qubit state / Bloch sphere &
Two-component complex vector algebra (\(\mathbb{CP}^1\cong S^2\));
non-unitary Jones matrices excluded & \textbf{{[}★ exact{]}} \\
9 & Phase noise / jitter & Dephasing (\(T_2\)) & Dissipation of the
phase distribution (Lindblad form on the quantum side, not decoherence
as a whole) & \textbf{{[}△ parallel{]}} \\
10 & SVD (linear-algebra tool) & Schmidt decomposition (bipartite pure
state) & Canonical form on a tensor-product space (capacity formulas not
isomorphic) & \textbf{{[}★ exact{]}} \\
\end{longtable}
}

\textbf{Exact identity}: 5 items (\#1, \#2, \#3, \#8, \#10)
\textbf{Strong structural correspondence}: 1 item (\#4)
\textbf{Structural parallel}: 4 items (\#5, \#6, \#7, \#9)

\begin{center}\rule{0.5\linewidth}{0.5pt}\end{center}

\subsection{Appendix B: Position of This
Note}\label{appendix-b-position-of-this-note}

This note is a standalone observation paper and has no aim of
reinforcing or extending any particular research program or body of
papers. All its references are classical original papers or standard
textbooks of the two fields, and it has no direct citation relation to
other work by the author.

We hope this note functions, for experts, educators, and
interdisciplinary researchers in the two fields, as a reference for the
mutual translation of concepts.

\begin{center}\rule{0.5\linewidth}{0.5pt}\end{center}

Author: Noriaki Kihara / WF System Co., Ltd.~/ ORCID
\href{https://orcid.org/0009-0004-6753-4020}{0009-0004-6753-4020} / CC
BY 4.0

\end{document}
